\clearpage
\SourceLeaf{328}
\pdfbookmark[1]{Remarks on the Supplements}{zusaetze-bemerkungen}
\begin{center}
\textbf{REMARKS.}
\end{center}

Notes [2], [4], [5], [13], [15] are found in a notebook, [3], [9], [10], [17], [18] on loose sheets, [1] and [12] in \textsc{Gauss}'s copies of the Berliner Astronomisches Jahrbuch for 1805 and 1818, and [6] in \textsc{Gauss}'s copy of \textsc{Vega}, log.-trig. Handbuch, Leipzig, 1797. Notes [7], [8], [11], [16] have already been published in the places indicated. Note [14] is taken from a notebook entitled »Abhandlungen von \textsc{Carl Friedrich Gauss}, Erstes Heft«, which was kindly made available by Professor \textsc{Jacobi} in Göttingen and comes from the papers of \textsc{Wilhelm Weber}. The name inscription »\textsc{J. H. Westphal}, Danzig 1818« is found inside the binding. In addition to some copies of papers by \textsc{Gauss}, the notebook contains notes of the same kind as Note [14], which appear to derive from lectures or other oral communications by \textsc{Gauss}.

The second part of Note [6] has also been included, although it has no direct connection with Articles 88---90. Compare Vol. VIII, p. 128.

The formulas in Note [15] are already found in the Disquisitio de elementis ellipticis Palladis (Vol. VI, pp. 17---18). Here they are put into a form more suited to numerical calculation. The result of the numerical example is also found there on p. 19. The discrepancy in the value of \(C\) arises because an arithmetic error in Note [15] was corrected for publication.

The loose sheets in the papers contain further individual derivations and calculations concerning the limits of the geocentric positions of celestial bodies (zodiac, Notes [16]---[18]). Preparing them for publication proved impracticable because they were too fragmentary.

Minor changes to the original text had to be made in several places.
\begin{flushright}
\textsc{Brendel}.
\end{flushright}
\begin{center}
\rule{0.3\linewidth}{0.4pt}
\end{center}
